\begin{definitionbox}{Definition (Homomorphism)}
\begin{definition}[Homomorphism]
\label{def:homomorphism}
A map $f$ is a \emph{homomorphism} for an operation $O$ when
\[
f(x\mathbin{O}y)=f(x)\mathbin{O'}f(y).
\]
For ring-like structures it preserves both addition and multiplication, and
preserves $1$ when the language includes a distinguished multiplicative
identity.
\end{definition}
\end{definitionbox}

\begin{remark*}[Standard quantified statement]
\[
\forall x,y\in A,\ f(x\mathbin{O}y)=f(x)\mathbin{O'}f(y).
\]
\end{remark*}

\begin{remark*}[Predicate reading]
\[
\operatorname{Homomorphism}(f,(A,O),(B,O')).
\]
\end{remark*}

\begin{remark*}[Interpretation]
A homomorphism preserves the named operation from the source structure to the
target structure.
\end{remark*}

\begin{dependencies}
\begin{itemize}
  \item \hyperref[def:function]{Function}
  \item \hyperref[def:binary-arithmetic-operation]{Binary Arithmetic Operation}
\end{itemize}
\end{dependencies}

\begin{definitionbox}{Definition (Injective Homomorphism)}
\begin{definition}[Injective Homomorphism]
\label{def:injective-homomorphism}
An \emph{injective homomorphism} is a homomorphism whose underlying function
is injective.
\end{definition}
\end{definitionbox}

\begin{remark*}[Standard quantified statement]
\[
\operatorname{Homomorphism}(f,(A,O),(B,O'))
\quad\text{and}\quad
\operatorname{Injective}(f,A,B).
\]
\end{remark*}

\begin{remark*}[Predicate reading]
\[
\operatorname{InjectiveHomomorphism}(f,(A,O),(B,O')).
\]
\end{remark*}

\begin{remark*}[Interpretation]
An injective homomorphism preserves the operation and identifies the source
with its image without collisions.
\end{remark*}

\begin{dependencies}
\begin{itemize}
  \item \hyperref[def:homomorphism]{Homomorphism}
  \item \hyperref[def:injective]{Injective Function}
\end{itemize}
\end{dependencies}

\begin{definitionbox}{Definition (Order Embedding)}
\begin{definition}[Order Embedding]
\label{def:abstract-order-embedding}
\[
xRy \Longleftrightarrow f(x)R'f(y).
\]
\end{definition}
\end{definitionbox}

\begin{remark*}[Standard quantified statement]
\[
\forall x,y\in A,\ xRy \Longleftrightarrow f(x)R'f(y).
\]
\end{remark*}

\begin{remark*}[Predicate reading]
\[
\operatorname{OrderEmbedding}(f,(A,R),(B,R')).
\]
\end{remark*}

\begin{remark*}[Interpretation]
An order embedding both preserves and reflects the relation.
\end{remark*}

\begin{dependencies}
\begin{itemize}
  \item \hyperref[def:function]{Function}
  \item \hyperref[def:relation]{Relation}
  \item \hyperref[def:order-embedding]{Order Embedding}
\end{itemize}
\end{dependencies}

\begin{definitionbox}{Definition (Structure Embedding)}
\begin{definition}[Structure Embedding]
\label{def:structure-embedding}
A \emph{structure embedding} is an injective homomorphism that also embeds the
relevant order relation.
\end{definition}
\end{definitionbox}

\begin{remark*}[Standard quantified statement]
\[
\operatorname{InjectiveHomomorphism}(f,(A,O),(B,O'))
\quad\text{and}\quad
\operatorname{OrderEmbedding}(f,(A,R),(B,R')).
\]
\end{remark*}

\begin{remark*}[Predicate reading]
\[
\operatorname{StructureEmbedding}(f,A,B).
\]
\end{remark*}

\begin{remark*}[Interpretation]
A structure embedding preserves the operation data and the relevant order data.
\end{remark*}

\begin{dependencies}
\begin{itemize}
  \item \hyperref[def:injective-homomorphism]{Injective Homomorphism}
  \item \hyperref[def:abstract-order-embedding]{Order Embedding}
\end{itemize}
\end{dependencies}

\begin{propositionbox}{Proposition (Composition of Homomorphisms)}
\begin{proposition}[Composition of Homomorphisms]
\label{prop:composition-of-homomorphisms}
The composition of two homomorphisms is a homomorphism.
\hyperref[prf:composition-of-homomorphisms]{\textit{Go to proof.}}
\end{proposition}
\end{propositionbox}

\begin{remark*}[Standard quantified statement]
\[
\operatorname{Homomorphism}(f,A,B)
\land
\operatorname{Homomorphism}(g,B,C)
\Longrightarrow
\operatorname{Homomorphism}(g\circ f,A,C).
\]
\end{remark*}

\begin{remark*}[Predicate reading]
\[
\operatorname{Homomorphism}(g\circ f,A,C).
\]
\end{remark*}

\begin{remark*}[Interpretation]
Operation preservation is stable under composition.
\end{remark*}

\begin{dependencies}
\begin{itemize}
  \item \hyperref[def:homomorphism]{Homomorphism}
\end{itemize}
\end{dependencies}

\begin{corollarybox}{Corollary (Identity Map Is a Homomorphism)}
\begin{corollary}[Identity Map Is a Homomorphism]
\label{cor:identity-map-is-homomorphism}
The identity map on a structure is a homomorphism from that structure to
itself.
\hyperref[prf:identity-map-is-homomorphism]{\textit{Go to proof.}}
\end{corollary}
\end{corollarybox}

\begin{remark*}[Standard quantified statement]
\[
\operatorname{Homomorphism}(\operatorname{id}_A,A,A).
\]
\end{remark*}

\begin{remark*}[Predicate reading]
\[
\operatorname{IdentityHomomorphism}(\operatorname{id}_A,A).
\]
\end{remark*}

\begin{remark*}[Interpretation]
The identity map preserves every operation of its own structure.
\end{remark*}

\begin{dependencies}
\begin{itemize}
  \item \hyperref[def:homomorphism]{Homomorphism}
\end{itemize}
\end{dependencies}
